%% ============================================================================
%% Seebeck-Prime — SoftwareX Original Software Publication manuscript
%%
%% Produced from the OFFICIAL Elsevier SoftwareX OSP template
%%   (templates/softwarex-osp-template.tex, downloaded unmodified on
%%   2026-11-12 from https://legacyfileshare.elsevier.com/promis_misc/
%%   softwarex-osp-template.tex).
%% The template's "Before you begin" instruction block was deleted as that
%% block itself instructs ("DELETE THIS SECTION BEFORE SUBMISSION"); its
%% mandatory requirements are recorded here as comments:
%%   * We only accept GitHub repositories within Metadata Tables (mandatory;
%%     README.md and Licence.txt required in the repository).
%%   * Word limit 4,000 words; max 6 figures; 5 mandatory templated sections.
%%   * Abstract "ca. 100 words" (guide maximum 250 words).
%%
%% Status: final assembly draft — NOT SUBMISSION-READY. Mandatory fields
%% marked [AUTHOR TO CONFIRM] must be resolved before submission (see
%% AUTHOR_INPUT_REQUIRED.md, AUTHOR_INFORMATION.md, SOFTWARE_METADATA.md,
%% REFERENCE_VERIFICATION.md).
%% ============================================================================

\documentclass[preprint,12pt, a4paper]{elsarticle}

\usepackage{amssymb}
\usepackage{hyperref}
\setlength{\parindent}{0pt}

\journal{SoftwareX}

\begin{document}

\renewcommand{\labelenumii}{\arabic{enumi}.\arabic{enumii}}

\begin{frontmatter}

%% ---------------------------------------------------------------------------
%% Title, authors, addresses
%% ---------------------------------------------------------------------------

\title{An integrated desktop application for thermoelectric
transport-property analysis and modelling (Seebeck-Prime:
reproducible Lorenz-conditioned analysis of thermoelectric data)}

%% [AUTHOR TO CONFIRM: exact title form. The SoftwareX convention is
%% "Title (SoftwareName: Then a short title)".]

\author{[AUTHOR TO CONFIRM: author name(s)]\corref{cor1}}
\cortext[cor1]{Corresponding author.}
\ead{[AUTHOR TO CONFIRM: corresponding author e-mail]}
\address{[AUTHOR TO CONFIRM: affiliation(s) and address(es)]}

%% ---------------------------------------------------------------------------
%% Abstract (template guidance: ca. 100 words; guide maximum 250)
%% ---------------------------------------------------------------------------

\begin{abstract}
Seebeck-Prime is an open-source desktop application and Python package for
thermoelectric transport-property analysis. It executes the complete chain from
measured Seebeck coefficient, electrical conductivity and thermal conductivity
to power factor, Lorenz-number-conditioned electronic and lattice thermal
conductivities, figure of merit and device efficiency as one integrated,
unit-normalised workflow, provided through a Qt interface for interactive use
and a headless API for reproducible scripting. Implemented models include three
Lorenz-number descriptions (constant, empirical, and rigorous
single-parabolic-band), Fermi--Dirac integrals, carrier density and
effective-mass estimation, Hall analysis, weighted mobility, laser-flash
thermal conductivity and first-order uncertainty propagation. Calculations are
verified against independently written reference implementations (relative
errors $10^{-14}$--$10^{-6}$) and against four published experimental datasets
(power factor and figure of merit within $1.3$--$5.9\,\%$). The software
explicitly separates numerical verification from physical validation and
documents its model assumptions and unverified comparisons.
\end{abstract}

%% ---------------------------------------------------------------------------
%% Keywords (guide: 1--7 keywords; draft trimmed from 10 to 7 for compliance)
%% ---------------------------------------------------------------------------

\begin{keyword}
thermoelectric materials \sep
Seebeck coefficient \sep
power factor \sep
Lorenz number \sep
single parabolic band \sep
Fermi-Dirac integrals \sep
figure of merit
\end{keyword}

\end{frontmatter}

%% ============================================================================
%% REQUIRED METADATA
%% ============================================================================

\section*{Required Metadata}
\label{}

\section*{Current code version}
\label{}

Ancillary data table required for subversion of the codebase. The left column
is retained verbatim from the template; the right column describes the current
state of Seebeck-Prime, with unresolved items marked \textbf{[AUTHOR TO
CONFIRM]}.

\begin{table}[!h]
\begin{tabular}{|l|p{6.5cm}|p{6.5cm}|}
\hline
\textbf{Nr.} & \textbf{Code metadata description} & \textbf{Please fill in this column} \\
\hline
C1 & Current code version & 0.1.0 (released 2026-11-13; reconciled across the changelog, package metadata and \texttt{\_\_version\_\_}) \\
\hline
C2 & Permanent GitHub link to code/repository used for this code version & \url{https://github.com/mostafa-kadi/Seebeck-prime} \textbf{[AUTHOR TO CONFIRM: repository must be made public before submission]}; the repository contains a documented \texttt{README.md}, \texttt{LICENSE} and \texttt{Licence.txt} \\
\hline
C3 & Legal Code License & MIT Licence, \url{https://opensource.org/licenses/MIT}; licence text shipped as \texttt{LICENSE} / \texttt{Licence.txt} in the repository \\
\hline
C4 & Code versioning system used & git \\
\hline
C5 & Software code languages, tools, and services used & Python (3.12+); NumPy, SciPy, pandas, Matplotlib, scienceplots, PyQt6, qtawesome; pytest \\
\hline
C6 & Compilation requirements, operating environments \& dependencies & \texttt{pip install -r requirements.txt}; Windows (primary); Python 3.12/3.13 per \texttt{docs/SUPPORTED\_PYTHON.md} \\
\hline
C7 & If available Link to developer documentation/manual & \textbf{[AUTHOR TO CONFIRM: public documentation URL]} --- local documentation exists (\texttt{README.md}, \texttt{docs/}) \\
\hline
C8 & Support email for questions & \textbf{[AUTHOR TO CONFIRM]} \\
\hline
\end{tabular}
\caption{Code metadata (mandatory)}
\label{}
\end{table}

\textbf{Main text}

Maximum 6 pages in total (excluding metadata, tables, figures, references) with
a 4000-word limit. The body of this manuscript is within the 4000-word limit.

%% ============================================================================
%% 1. MOTIVATION AND SIGNIFICANCE
%% ============================================================================

\section{Motivation and significance}
\label{sec:motivation}

Thermoelectric transport-property analysis is a chain of interdependent
quantities. The measured Seebeck coefficient $S$, electrical conductivity
$\sigma$ and total thermal conductivity $\kappa$ are combined into the power
factor $PF = S^{2}\sigma$ and the dimensionless figure of merit
$ZT = PF\,T/\kappa$ \cite{snyder2008,goldsmid2016}. The total thermal
conductivity is then decomposed into electronic and lattice contributions,
$\kappa = \kappa_{e} + \kappa_{L}$, which requires a choice of Lorenz number
$L$ in the Wiedemann--Franz relation $\kappa_{e} = L\sigma T$.

This chain is where preventable error accumulates. The lattice contribution is
a \emph{difference of two comparable numbers}, so an error in the Lorenz number
propagates directly and substantially into $\kappa_{L}$ and hence into $ZT$.
The choice of Lorenz number is also where the physics becomes subtle: the
constant Sommerfeld value, an empirical relation $L(S)$ \cite{kim2015}, and a
rigorous single-parabolic-band (SPB) transport-integral model
\cite{maysnyder2012} each embed different assumptions about scattering and band
structure, and the differences are largest exactly where the subtraction is
most sensitive.

In practice, researchers assembling these steps from separate programs,
spreadsheets or custom scripts face inconsistent unit conventions and implicit
model choices. Seebeck-Prime was developed so that the chain --- from input,
through calculation, to derived properties, plots and exports --- is executed,
unit-normalised and traceable within one application. Its scientific problem
statement is: \emph{compute the thermoelectric transport chain reproducibly,
make every Lorenz-number assumption explicit, and refuse silent unit or model
errors.}

The intended experimental setting (how the user uses the software) is
twofold. For interactive analysis, the Qt desktop application supports manual
entry, instrument-file import, batch and multi-sample processing, plotting and
export. For reproducible work, the same calculations are exposed as a headless
Python API, and the validation programme itself is expressed as an executable
benchmark suite against published data \cite{katsura2025}.

Related work: the transport-model background follows the SPB modelling
conventions of May and Snyder \cite{maysnyder2012}, the Lorenz-number
estimation context follows Kim et al.\ \cite{kim2015}, and the weighted-mobility
formulation follows Snyder et al.\ \cite{snyder2020}. Benchmark source data are
drawn from published thermoelectric measurements digitised in the Starrydata
project \cite{katsura2025}.

%% ============================================================================
%% 2. SOFTWARE DESCRIPTION
%% ============================================================================

\section{Software description}
\label{sec:software}

\subsection{Software architecture}
\label{sec:architecture}

Seebeck-Prime follows a layered architecture
(Fig.~\ref{fig:architecture}):

\begin{itemize}
  \item \textbf{Presentation:} Qt desktop user interface (\texttt{src/gui}),
        plotting layers (\texttt{src/presentation}), reusable UI helpers;
  \item \textbf{Application:} use cases and orchestration
        (\texttt{src/application}): calculation workflows, validation
        workflows, file-handling orchestration;
  \item \textbf{Domain:} scientific models and pure calculation functions
        (\texttt{src/domain}): transport equations, SPB/Fermi-integral model,
        Hall, phonon, device efficiency, uncertainty, units, physical
        constants, validation rules;
  \item \textbf{Infrastructure:} storage and I/O adapters.
\end{itemize}

Architectural boundary tests enforce the direction of dependencies (no
domain$\to$GUI, no infrastructure$\to$GUI), keeping the scientific core free
of UI concerns.

\begin{figure}[htbp]
\centering
\fbox{\parbox[c][3.5cm][c]{0.9\textwidth}{\centering
\textbf{[Artwork pending]}\\
Layered-architecture diagram --- see \texttt{FIGURE\_PLAN.md} (Fig.~1).}}
\caption{Layered architecture of Seebeck-Prime: presentation, application,
domain and infrastructure layers, with enforced dependency boundaries.}
\label{fig:architecture}
\end{figure}

\begin{figure}[htbp]
\centering
\fbox{\parbox[c][3.5cm][c]{0.9\textwidth}{\centering
\textbf{[Screenshot pending --- author must capture]}\\
Genuine screenshot of the running application ---
see \texttt{FIGURE\_PLAN.md} (Fig.~2). No screenshot may be fabricated.}}
\caption{Main window of the Seebeck-Prime desktop application.}
\label{fig:gui}
\end{figure}

\subsection{Software functionalities}
\label{sec:functionalities}

The core headless entry point computes, from temperature, electrical
conductivity and Seebeck coefficient (and total thermal conductivity), the
four central transport relations:

\begin{equation}
PF = S^{2}\,\sigma ,
\end{equation}
\begin{equation}
\kappa_{e} = L\,\sigma\,T ,
\end{equation}
\begin{equation}
\kappa_{L} = \kappa - \kappa_{e} ,
\end{equation}
\begin{equation}
ZT = \frac{PF\,T}{\kappa} .
\end{equation}

All quantities are handled internally in SI base units and converted for
display, with unit validation on conversion (a thermal-conductivity
unit-convention defect, severity SE-042, was identified and remediated during
the validation programme).

Three Lorenz-number modes are implemented:

\begin{itemize}
  \item \textbf{Constant:} $L_{0} = (\pi^{2}/3)(k_{B}/e)^{2} =
        2.443004509\times10^{-8}\ \mathrm{W\,\Omega\,K^{-2}}$; the
        application-wide default is $2.44\times10^{-8}\
        \mathrm{W\,\Omega\,K^{-2}}$;
  \item \textbf{Kim empirical:} $L(S) = (1.5 +
        e^{-|S|/116})\times10^{-8}\ \mathrm{W\,\Omega\,K^{-2}}$
        \cite{kim2015};
  \item \textbf{Rigorous SPB:} $L(\eta)$ via Fermi--Dirac transport
        integrals \cite{maysnyder2012}.
\end{itemize}

Further implemented functionality, at the headless level, includes:
SPB Fermi--Dirac integrals, Seebeck-from-reduced-Fermi-level inversion,
carrier concentration and effective-mass estimation, and round trips; Hall
coefficient/density/mobility and scattering-exponent fit; weighted mobility
\cite{snyder2020}; laser-flash thermal conductivity
$\kappa = \alpha\,C_{p}\,\rho$, with laboratory inputs converted to SI before
evaluation; device efficiency; first-order uncertainty propagation of power
factor, thermal-conductivity components and $ZT$ (input independence is
assumed); unit conversion; manual and instrument-derived workflows, batch and
multi-sample processing; GUI, plotting and export.

A lattice-thermal-conductivity diagnostic is wired into the GUI and
multi-sample reports: when the subtraction $\kappa_{L} = \kappa - \kappa_{e}$
yields a negative value, the software surfaces a persistent, explicit notice
(the value itself is never silently clipped), because a negative $\kappa_{L}$
is physically impossible and almost always indicates a Lorenz-model or input
inconsistency.

%% ============================================================================
%% 3. ILLUSTRATIVE EXAMPLES
%% ============================================================================

\section{Illustrative examples}
\label{sec:examples}

\subsection{Headless calculation of the transport chain}
\label{sec:example-headless}

The same equations are available programmatically. The following example uses
the public headless API (SI units throughout):

\begin{verbatim}
from seebeck_prime.domain.calculations import (
    calculate_pf, calculate_kappa_e, calculate_kappa_l, calculate_zt,
)
from seebeck_prime.domain.physical_constants import DEFAULT_LORENZ

pf = calculate_pf(seebeck, electrical_conductivity)        # S [V/K], sigma [S/m]
kappa_e = calculate_kappa_e(electrical_conductivity, temperature,
                            lorenz=DEFAULT_LORENZ)          # L [W*Ohm/K^2]
kappa_l = calculate_kappa_l(kappa_total, kappa_e)
zt = calculate_zt(pf, temperature, kappa_total)
\end{verbatim}

\noindent
The functions return $PF$ in W/(m\,K$^{2}$), $\kappa_{e}$ and $\kappa_{L}$ in
W/(m\,K), and dimensionless $ZT$; the default Lorenz argument is the
application-wide constant $2.44\times10^{-8}\ \mathrm{W\,\Omega\,K^{-2}}$.

\subsection{Validation against published experimental data}
\label{sec:example-benchmarks}

The executable benchmark suite compares the calculation engine against four
published specimens whose raw curves were extracted from the Starrydata
project \cite{katsura2025} with checksum verification: an n-type
Bi$_{2}$Te$_{2.4}$Se$_{0.6}$ alloy \cite{huang2019}, a pseudo-ternary
(PbTe)$_{0.95-x}$(Sb$_{2}$Se$_{3}$)$_{x}$(PbS)$_{0.05}$ specimen
\cite{zhang2020}, an Sb-doped Si$_{80}$Ge$_{20-x}$Sb$_{x}$ multiphase alloy
\cite{wang2015}, and the single-phase P-doped Si$_{80}$Ge$_{20}$ reference of
an SrTiO$_{3}$/SiGe composite study \cite{wang2020srtio3}. Documented
outcomes:

\begin{itemize}
  \item Published power factors and $ZT$ are reproduced within approximately
        $1.3\,\%$--$5.9\,\%$ across the four specimens.
  \item Published Si$_{80}$Ge$_{20}$P lattice thermal conductivity is
        reproduced within $0.7\,\%$ using the application Lorenz constant.
  \item The Lorenz number implied by the published
        Si$_{80}$Ge$_{20}$P data ($2.46\times10^{-8}\
        \mathrm{W\,\Omega\,K^{-2}}$) is consistent with the constant
        ($2.44\times10^{-8}$).
  \item The acoustic-scattering SPB Lorenz value ($1.59\times10^{-8}$)
        disagrees by approximately $35\,\%$; investigation showed this to be a
        scattering-assumption matter (effective scattering exponent
        $\lambda \approx 1.41$), not an implementation error, and the model was
        deliberately \emph{not} refitted.
  \item The PbTe comparison is reported \textbf{unverified} because the
        published $\kappa/\kappa_{L}$ pair is internally inconsistent (spread
        $>20\,\%$).
\end{itemize}

These outcomes are stated with their honest boundaries: tests passing is not
validation; numerical agreement is not experimental validation; conditional
physical models are not universal.

%% ============================================================================
%% 4. IMPACT
%% ============================================================================

\section{Impact}
\label{sec:impact}

No user, download or adoption statistics are claimed; the software has not yet
been publicly released (see \textbf{Current executable software version}). The
expected impact rests on the demonstrated capability:

\begin{itemize}
  \item \textbf{Reproducible analysis chain.} Researchers obtain
        $PF$, $\kappa_{e}$, $\kappa_{L}$, $ZT$ and device efficiencies from
        measured inputs through one unit-normalised path, in both interactive
        and scripted form.
  \item \textbf{Explicit Lorenz-model sensitivity.} Because
        $\kappa_{L} = \kappa - L\sigma T$ is a difference of comparable
        numbers, the Lorenz choice is the dominant avoidable error in
        $\kappa_{L}$ and $ZT$. Seebeck-Prime makes each Lorenz assumption
        explicit, offers constant/empirical/SPB models side by side, and
        surfaces physically impossible results (negative $\kappa_{L}$) instead
        of hiding them. The benchmark finding that a published
        Si$_{80}$Ge$_{20}$P dataset implies the Sommerfeld value while the
        acoustic SPB value disagrees by $\sim 35\,\%$ (scattering exponent
        $\lambda \approx 1.41$) is exactly the kind of model-sensitivity
        analysis the software is designed to make routine.
  \item \textbf{Reusable benchmarking.} The checksum-verified benchmark
        framework against Starrydata-drawn datasets \cite{katsura2025} can be
        pointed at further published data, including cases where the honest
        verdict is ``unverified'' (as for the PbTe specimen).
  \item \textbf{Research questions enabled.} How do Lorenz-model choices and
        scattering assumptions propagate into reported lattice thermal
        conductivities and figures of merit across digitised datasets; and how
        should analysis software communicate model conditionality to users.
\end{itemize}

%% ============================================================================
%% 5. LIMITATIONS AND FUTURE WORK
%% ============================================================================

\section{Limitations and future work}
\label{sec:limitations}

Documented limitations, stated so that no capability is overstated:

\begin{itemize}
  \item Automated GUI test coverage is incomplete (the non-modal dialogs lack
        systematic test coverage).
  \item A raw voltage/current $\to$ resistance $\to$ resistivity $\to$
        conductivity chain is \textbf{not implemented}.
  \item Covariance-aware uncertainty propagation ($\mathbf{J}\,\Sigma\,\mathbf{J}^{T}$)
        is \textbf{not implemented}; current propagation assumes independent
        inputs.
  \item Kane/Dingle/multiband Lorenz models are \textbf{not implemented}.
  \item The weighted-mobility implementation is \textbf{not experimentally
        validated}; its behaviour is documented, including documented
        divergence from an earlier internal formulation.
  \item The PbTe and weighted-mobility benchmark comparisons remain
        \textbf{unverified}.
  \item Public release, licensing finalisation and version reconciliation are
        pending (release-phase items).
\end{itemize}

Future work prioritises: a public GitHub release with licence and archival;
covariance-aware uncertainty propagation; the raw V/I measurement chain as a
separate modular workflow; and systematic GUI test coverage.

%% ============================================================================
%% 6. CONCLUSIONS
%% ============================================================================

\section{Conclusions}
\label{sec:conclusions}

Seebeck-Prime implements the integrated thermoelectric transport chain ---
power factor, Lorenz-conditioned electronic and lattice thermal conductivities,
figure of merit and device efficiency --- with traceable units, explicit
Lorenz-model alternatives, and a headless API equivalent to the GUI workflow.
Its calculations are numerically verified against independently written
reference implementations to relative errors of $10^{-14}$--$10^{-6}$, and its
benchmarks against four published experimental specimens reproduce published
power factors and figures of merit within $1.3$--$5.9\,\%$, with the honest
boundaries (unverified comparisons, unvalidated models, unimplemented
extensions) documented rather than hidden.

%% ============================================================================
%% DECLARATIONS
%% ============================================================================

\section*{Declaration of competing interest}
\textbf{[AUTHOR TO CONFIRM: competing-interest declaration --- the standard
form is ``The authors declare that they have no known competing financial
interests or personal relationships that could have appeared to influence the
work reported in this paper.'']}

\section*{Acknowledgements}
\textbf{[AUTHOR TO CONFIRM: acknowledgements, funding, or ``none'']}

%% ---------------------------------------------------------------------------
%% Generative-AI disclosure (required before the reference list)
%% ---------------------------------------------------------------------------

\section*{Declaration of generative AI and AI-assisted technologies in the
writing process}
\textbf{Statement:} During the preparation of this work the author(s) used
\textbf{[AUTHOR TO CONFIRM: AI tool name(s)/version(s)]} in order to draft,
organise and check the internal consistency of the manuscript, based on
verified repository evidence. After using this tool, the author(s) reviewed
and edited the content as needed and take(s) full responsibility for the
content of the published article. See \texttt{AI\_USAGE\_RECORD.md} for the
complete record.

%% ============================================================================
%% REFERENCES
%% ============================================================================

\bibliographystyle{elsarticle-num}
\bibliography{references}

%% ============================================================================
%% CURRENT EXECUTABLE SOFTWARE VERSION
%% ============================================================================

\section*{Current executable software version}
\label{}

\begin{table}[!h]
\begin{tabular}{|l|p{10.5cm}|}
\hline
\textbf{Item} & \textbf{Value} \\
\hline
Permanent link & \url{https://github.com/mostafa-kadi/Seebeck-prime} \textbf{[AUTHOR TO CONFIRM: repository must be made public before submission]} \\
\hline
Legal Code License & MIT Licence (\texttt{LICENSE} / \texttt{Licence.txt} in the repository) \\
\hline
Current code version & 0.1.0 \\
\hline
Current version date & 2026-11-13 \\
\hline
Code repository & \url{https://github.com/mostafa-kadi/Seebeck-prime} \textbf{[AUTHOR TO CONFIRM: repository must be made public before submission]} \\
\hline
Related research manuscripts & This article. \\
\hline
\end{tabular}
\caption{Current executable software version}
\label{}
\end{table}

\end{document}
